\chapter{Computational and Reproducibility Record}\label{app:reproducibility}

\section{Completed numerical records}
The baseline results are associated with the computational record
\path{basic_particle_problem/results/final_study}. The Gaussian
verification results and the diagnostic coupled calculations are associated
with \path{final_cloud_code/outputs/evaluated_2d_final}. These two result
sets are kept separate from short environment smoke tests. The smoke tests
are not used as research results in this thesis.

The archived repository revisions are
\begin{quote}\small\ttfamily
basic\_particle\_problem:\par
c3eb0efe76bcd62a3299dcd6c28922dbfa4c667a\par\medskip
final\_cloud\_code:\par
9d996c4f555bf5bf09e9f6fa78a3877e44310ec9
\end{quote}
The source bundle contains a SHA-256 manifest of the numerical inputs and
method files used to prepare the manuscript. The original result files are
copied without numerical modification. Figure scripts change presentation,
not the stored solution values. Gauge alignment and the exact computations
used to recompute metrics are explicit in those scripts.

The historical result manifests retain their own source hashes. The
archived baseline model and geometry, and the Gaussian model, physics, and
verification modules match their recorded hashes. Current driver and
configuration files do not all match the historical manifests; they are
retained for method inspection and are not claimed as an exact replay of
the original training execution.

\section{Training and evaluation distinction}
The baseline final evaluation restores previously trained network weights.
The associated original flow configuration uses 1,000 Adam epochs, a
250-iteration L-BFGS budget, and a refinement budget of 1,800 epochs with
700 L-BFGS iterations. The original force budget is 1,000 epochs. The two
flow seeds are 42 and 43. The later evaluation configuration contains
different default training fields, which do not imply that restored weights
were retrained with those defaults.

The Gaussian channel uses seed 7, float64, and the evaluated configuration
with 1,200 fluid/parcel Adam steps and a 200-iteration L-BFGS budget. Its
reduced ODE tests use 700 Adam steps with a 120-iteration refinement budget;
the Taylor--Green test uses 1,200 Adam steps with a 200-iteration refinement
budget. An iteration budget need not equal the number of closure
evaluations in a line-search optimizer. The stored status files and
configuration distinguish these quantities.

Reference values are excluded from the Gaussian reduced-test optimization.
In V01 the initial response is imposed. In V02 the initial values and total
momentum relation are imposed. In V10 the initial velocity and pressure
structure are imposed while the temporal decay is held out. These choices
determine which aspects of the tests are independent predictions.

\section{Software and execution scope}
The numerical environment is Linux with Python 3.12. The baseline CPU
environment uses PyTorch 2.14.1, and the Gaussian CPU environment uses
PyTorch 2.6.0. The scientific stack includes NumPy 2.3.5, SciPy 1.17.0,
Matplotlib 3.10.8, and scikit-fem 12.0.2. The preserved historical Gaussian
runtime file records CPU execution with two threads. These are
computational reproducibility details; no GPU performance or speedup is
claimed.

The source bundle supports two levels of reproduction. Compiling
\path{thesis.tex} uses the included tables and vector figures and does not
require retraining. Rebuilding the figures additionally uses the archived
arrays, model definitions, and two small baseline checkpoints. Independent
metric checks include the full Poiseuille velocity vector, centered
pressure, the final held-out drag evaluation, and the reduced Gaussian
benchmarks.

Reproducing all training runs requires the complete corresponding
repositories and their dependency environments. It also requires their
reference-input files. The manuscript bundle is a reproducible evidence
and typesetting package, not a claim that every full simulation can be
retrained from the manuscript alone.

\section{Figures and numerical transformations}
Nineteen figures are included, each registered with its input source or
its schematic status. Architecture figures are explanatory drawings.
The constriction and kernel-wall illustrations specify design choices.
They are never registered as computed flow results. Numerical plots retain
all evaluation arrays for aggregate errors; marker subsampling is used only
to make reference curves readable.

Pressure comparisons subtract an additive constant independently from each
array. Paired field panels use common color limits. Discrepancy panels have
their own explicitly described scale and are not smoothed to conceal local
error. The drag parity plot uses a stored diagnostic set different from
the final held-out metric set; both are identified in the caption and
validation record.

Five additional figures document the finite coupled parcel response of the
identified $\St=0.6$, $\beta=0.05$ calculation. Their diagnostic status and
the unresolved fluid residual are explicit in \cref{ch:gaussian}; they
are not relabeled as accepted predictive evidence. No synthetic
demonstration or extruded three-dimensional illustration is used as a
simulation result.
The material limitations are retained in the main chapters and the complete
baseline trajectory table. Additional unsuccessful run histories remain
outside the figure selection because they do not add accepted transport
results.

